% \kenneth{TODO TY: If possible, replace the footnotes with citations. Footnotes eat space.}

A total of 14 methods were included for comparison: 
six reuse baselines with six input variations
% (Mention, Paragraph, Window[0, 1], Window[0, 2], Window[1, 1], and Window[2, 2]); 
(M, P, W[0, 1], W[0, 2], W[1, 1], and W[2, 2]);
five text summarization models with five inputs
% (Mention, Mention+OCR, Paragraph, Paragraph+OCR, and OCR); 
(M, M+O, P, P+O, and O);
one text summarization model using P+O with controlled data quality; 
and two vision-to-language models (BEiT+GPT-2 and TrOCR). 
Note that we use subscripts of M, P, W, O, B to denote the input features: 
Mention, Paragraph, Window, OCR, and Better data quality, respectively.
% For example, Reuse$_{W[0, 1]}$ refers to the Reuse baseline outputting Window[0, 1] and Pegasus$_{P+O+B}$ refers to the text summarization model using better quality of Paragraph+OCR as the input.
% \kenneth{We could probably remove FOR EXAMPLE to save space.}
The model training details and decoding configuration are provided in \Cref{sec:appendix-model-training-details}.
% Next, we describe the model training details for both the text summarization models and the vision-to-language models.
% Reuse$_{M}$, Reuse$_{P}$, Reuse$_{W[0, 1]}$
% Pegasus$_{M}$, Pegasus$_{M+O}$, Pegasus$_{M+O+B}$

% \paragraph{Model Training Detail.}\label{sec:appendix-model-training}
% \cy{Need update}

% \kenneth{TODO CY: Move the following details to Appendix.}

%We then applied both automatic evaluation and human evaluation to compare model performance.


